% Generated draft: requires mathematical and source review.
\documentclass[11pt]{article}
\usepackage[margin=25mm]{geometry}
\usepackage{fontspec}
\setmainfont{Latin Modern Roman}
\usepackage{amsmath,amssymb,mathtools,unicode-math}
\setmathfont{Latin Modern Math}
\usepackage{enumitem,needspace,xurl,hyperref}
\setlength{\parindent}{0pt}
\setlength{\parskip}{5pt}
\setlength{\emergencystretch}{3em}
\newcommand{\N}{\mathbb{N}}
\newcommand{\Z}{\mathbb{Z}}
\newcommand{\Q}{\mathbb{Q}}
\newcommand{\R}{\mathbb{R}}
\newcommand{\C}{\mathbb{C}}
\newcommand{\F}{\mathbb{F}}
\newcommand{\A}{\mathbb{A}}
\newcommand{\PP}{\mathbb{P}}
\newcommand{\E}{\mathbb{E}}
\newcommand{\eps}{\varepsilon}
\newcommand{\dd}{\,\mathrm{d}}
\begin{document}
\section*{UnsolvedMath ID: 700101}
\textbf{Source problem code:} catalogue entry
\subsection{Definitions and mathematical statement}
Let \(G\) be a group. A strict total order on \(G\) is an irreflexive, transitive relation \(<\) such that, for distinct \(x,y\in G\), exactly one of \(x<y\) and \(y<x\) holds. Such an order is left invariant if
\[
\forall g,x,y\in G,\qquad x<y\Longrightarrow gx<gy,
\]
and two-sided invariant if
\[
\forall a,b,x,y\in G,\qquad x<y\Longrightarrow axb<ayb.
\]
Write \(\mathcal O_L(G)\) and \(\mathcal O_B(G)\) for the sets of left-invariant and two-sided invariant strict total orders, respectively. Distinct orders are counted as distinct relations on the underlying group, without identifying them under automorphisms.

The supplied record does not define the compatibility intended by ``linearly orderable.'' Consequently the following formulations are conditional on that convention: take \(\sigma=B\) if two-sided invariance is intended, and \(\sigma=L\) if left invariance is intended. These are alternative readings, not two separately attributed problems.

\begin{enumerate}
\item[(a)] \textup{(Attributed to L. Fuchs.)} Describe, up to group isomorphism, all groups \(G\) such that
\[
0<\lvert\mathcal O_\sigma(G)\rvert<\aleph_0.
\]
The positive lower bound expresses the requirement that the group be linearly orderable.
\item[(b)] \textup{(Supplied archived solved part.)} Does there exist a group \(G\) such that
\[
\lvert\mathcal O_\sigma(G)\rvert=\aleph_0?
\]
This uses ``countably many'' to mean countably infinitely many. If the source instead means at most countably many compatible orders, the corresponding condition for an orderable group is \(0<\lvert\mathcal O_\sigma(G)\rvert\leq\aleph_0\). The supplied excerpt does not explicitly distinguish these meanings.
\end{enumerate}

The record labels the overall problem partially solved, retains part (a) in the main section, and reports an affirmative answer to part (b), attributed to R. N. Buttsworth. This is a report of the supplied archive assessment; it does not verify that answer or resolve the ordering convention.
\subsection{Short English statement}
Part (a), attributed to L. Fuchs, asks for a classification of groups admitting a positive finite number of compatible linear orders. Part (b) asks whether a group can admit countably infinitely many such orders; the supplied archive reports an affirmative answer. The intended invariance of the orders and the precise meaning of ``countably many'' require source clarification.
\subsection{Sources}
\begin{itemize}
\item Local UnsolvedMath registry, permanent ID \(700101\), source problem code \texttt{catalogue entry}, title \emph{Kourovka Notebook Problem 2.25}; supplied complete source record.
\item \emph{The Kourovka Notebook (Unsolved Problems in Group Theory)}, 2nd issue (1967), Problem 2.25; proposed by A. I. Kokorin. Part (a) is separately attributed in the statement to L. Fuchs. The supplied background identifies an update in the 21st issue (2026), \texttt{arXiv:1401.0300v48}; the collection description names E. I. Khukhro and V. D. Mazurov as editors.
\item Supplied Archive of Solved Problems entry for part (b): affirmative answer attributed to R. N. Buttsworth, \emph{Bull. Austral. Math. Soc.} \textbf{4} (1971), 97--104. No article title or external link is supplied.
\item Supplied status review dated \texttt{2026-10-16}: canonical status \texttt{partially\_solved}, evidence field \texttt{published}; solved part(s) are recorded in the archive and remaining part(s) in the main section. This assessment is supplied dated curation, not independently verified source or literature checking.
\end{itemize}
\end{document}
